\section{Overlap removal}
\label{sec:ch6_overlap}

During the object reconstruction, two different object could share the same detector measurements.
An ordered overlap-removal procedure is therefore applied to resolve these ambiguities before event selection.

Geometrical overlaps are evaluated using $\Delta R_y$, obtained by replacing $\eta$ with rapidity $y=\tfrac{1}{2}\ln[(E+p_z)/(E-p_z)]$ in Eq.~\eqref{eq:delta_r}.
Muon--jet matches are also identified using \emph{ghost association}, in which a particle with negligible momentum is added along the muon direction during jet clustering.
The muon is associated with the jet containing this particle, without changing the jet momentum.

The following steps are applied in order to the baseline objects, using only candidates retained by the preceding steps:
\begin{enumerate}
    \item \textbf{Electron--electron}: the lower-$\pT$ electron is removed if two electrons share an ID track.
    \item \textbf{Tau--electron}: a tau candidate within $\Delta R_y<0.2$ of an electron is removed.
    \item \textbf{Tau--muon}: a tau candidate within $\Delta R_y<0.2$ of a muon is removed.
    \item \textbf{Calorimeter-tagged muon--electron}: a calorimeter-tagged muon sharing the electron ID track is removed.
    \item \textbf{Electron--muon}: an electron sharing an ID track with a remaining muon is removed.
    \item \textbf{Jet--electron}: a jet within $\Delta R_y<0.2$ of an electron is removed, retaining the electron interpretation of the EM shower.
    \item \textbf{Electron--jet}: an electron within $\Delta R_y<0.4$ of a remaining jet is removed to reject electrons embedded in hadronic jets.
    \item \textbf{Jet--muon}: a jet with fewer than three associated tracks is removed if the muon is ghost-associated with it. The small track multiplicity suggests that the jet candidate is dominated by activity associated with the muon.
    \item \textbf{Muon--jet}: a muon within $\Delta R_y<0.4$ of a remaining jet is removed.
    \item \textbf{Jet--tau}: a jet within $\Delta R_y<0.2$ of a remaining tau candidate is removed.
\end{enumerate}

